% Nouveau EGA V (cote 25): archive cote 25, site page 172 (PDF page 173).
% English translation of the checked French annotated master.
% AI-translated by GPT-6 Astra (OpenAI, ChatGPT Pro mode).
\scanpage{25/172}
\mspage{\hand{7}}

EGA IV 19

Second thoughts on ``Singular and supersingular set ...'' .

Having just written up the sorites of transversality, I realize that \unsure{the}
context adopted in my notes on singular and \unsure{supersingular} \unsure{sets},
namely sections of loc free Modules, i.e. good vector bundles,
is still not the ``right'' one. If $u$ is a section of a vector bundle
$E = W(\underline{E})$ (associated with $\underline{E}$ loc free of finite type) over X smooth over Y, then $V(u)$
is simply the inverse image \unsure{of the} \var{\inserted[typed]{subprescheme $E_{o}$ of $E$}}{subprescheme $E_{o}$ of $E$}{insertion}{ega25_25_172-v0004} ``zero section'', and a point $x\hand{\in} V(u)$
is nonsingular \var{\inserted[typed]{for u}}{for u}{insertion}{ega25_25_172-v0005} rel. to S iff u is \uline{transverse} to $E_{o}$ at x rel. to Y.
If this condition is not satisfied, i.e. if x is singular for u rel to
Y, we study the conditions under which x is not too singular
rel. to Y, i.e. u does not deviate too far from transversality at x. I have
no doubt that all the developments in my notes, with essentially
trivial changes, extend to the case where one has an X-prescheme E, a
subprescheme $E_{o}$ of E, a section u of E over X, and considers the inverse
image $X_{o}$ of $E_{o}$ under u, \var{\inserted[typed]{assuming E and $E_{o}$ smooth over X at the points of $v(X)$.}}{assuming E and $E_{o}$ smooth over X at the points of $v(X)$.}{insertion}{ega25_25_172-v0008} it is possible (but \unsure{only} a detailed
account will enable one to detect this) that it will sometimes be necessary to assume that
the relative dimension of $E_{o}$ over X at the points of $v(X_{o})$ is zero, i.e. that $E_{o}$ is
even étale over X at these points (which will be the case when $E_{o}$ comes, as
above, from a section of E over X, say $v_{o}$). Thus, x will be a point
``\var{\xout{singular}}{}{typewriter strike-out}{ega25_25_172-v0011} \var{\inserted[typed]{of nontransversality to $E_{o}$}}{of nontransversality to $E_{o}$}{insertion}{ega25_25_172-v0012} that is ordinary for v, relative to Y'' if and only if \var{\inserted[typed]{$x$}}{$x$}{insertion}{ega25_25_172-v0013} \var{\deleted{$\illegible{1}_{oy}$}}{}{deletion}{ega25_25_172-v0014} \var{\inserted{$x$}}{$x$}{insertion}{ega25_25_172-v0015} is
geometrically an ordinary quadratic point of $X_{o\,y}$ , and if $\mathrm{codim}_{x}(X_{oy},X_{y})$
$= c$ , where c is the codimension of $E_{o}$ in E (at $v(x)$) ; this implies that
$X_{o}\hand{\longrightarrow} X$ is \var{\inserted{an imm.}}{an imm.}{insertion}{ega25_25_172-v0016} transversely regular at x rel. to Y .... The role of the Module
$\underline{E}\hand{\otimes}\underline{O}_{V(u)}$ will here be played by $v_{o}^{\hand{*}}(\underline{N}_{E_{o}/E})$ , where $v_{o}:X_{o}\hand{\longrightarrow} E_{o}$ is induced by v. I
think you will have no trouble writing everything in this context, which is much more
natural after all---with the translations to the vector bundle case
left to the end. If there are difficulties even in getting started, we
could work on them \unsure{to}\unsure{gether} when you next pass through Paris ... \hand{I shall send you brief notes in the next few days.}
\var{\margin{\hand{172}}}{\placedmargin{\hand{172}}}{marginal note placed}{ega25_25_172-v0018}
\tnote{25-172-1}{The doubts in the, sets, supersingular and of the retain respectively the French e/r, h/l, final s and a/u uncertainties. Only retains the s/g doubt in seul; the two parts of together retain the e/t and b/v doubts in ensemble. The canceled s in the article les, the doubly canceled la, the isolated canceled f and the feminine e in régulière have no separate English effect. The canceled fragment d\unsure{e|u} remains doubtful. Within the inserted clause, the overwritten old \unsure{O|e} and \unsure{ant} likewise have no English spelling effect.}
\tnote{25-172-2}{Loc means locally; rel. means relative; sss is rendered iff; imm. is retained for immersion. The words that is allow ordinary to follow the inserted complement in English while point remains outside the insertion. The source first says S and then Y, changes from u to v, uses the two displayed occurrences of x in succession, and uses the indicated meanings of the subscripted v; these readings are retained.}
